\chapter{Limitaciones y Prospectiva}

Este capítulo afronta, con criterio académico, el alcance y los límites de lo desarrollado, y ofrece líneas de mejora o de investigación futura. Complementa las conclusiones, orientadas a cierre, con una visión crítica y hacia adelante.

\section{Limitaciones técnicas y de datos}
Enumérense las restricciones relevantes: supuestos de los modelos (por ejemplo, elección de arquitectura, duración o representatividad de la muestra), riesgo de \textit{overfitting}, sensibilidad a hiperparámetros, coste computacional, o incompatibilidad entre frecuencia de datos y horizonte de inversión. Cada limitación debería enlazarse con su posible impacto en los resultados, sin convertirse en mera lista.

\section{Limitaciones de interpretación o generalización}
Indique los contextos de mercado o de política (normativa, trato fiscal, acceso a productos) en que las conclusiones deben leerse con precaución, o el grado en que la experimentación es replicable en otros mercados o con otras clases de activo.

\section{Líneas de investigación y extensión}
Propóngase mejoras técnicas (nuevas arquitecturas generativas, regularización, criterios de riesgo alternativos, integración con restricciones normativas), ampliación de bases de datos, o estudios de robustez fuera de muestra. Cierre con un párrafo que ordene esas vías según plausibilidad o interés, sin obligar a una agenda rígida.
